\documentclass[../../../include/open-logic-section]{subfiles}

\begin{document}

\olfileid{his}{set}{mythology}

\olsection{इतिहासापेक्षा दंतकथाच अधिक?}

या घटनांकडे शतकाहून अधिक काळानंतर मागे वळून पाहताना, त्यांच्याबद्दल प्रचलित झालेले निष्कर्ष विनातपासणी स्वीकारू नयेत. ते गणिती निष्कर्ष नक्कीच नवे, उत्सुकता वाढवणारे आणि आश्चर्यकारक होते. पण ते खरोखर किती धक्कादायक होते? आणि भूमितीय अंतर्ज्ञानावर विसंबू नये, हे त्यांनी खरोखर दाखवले का?

धक्क्याच्या प्रश्नावर \citet{Gouvea2011} असे निदर्शनास आणतात की कँटर यांनी डेडेकिंड यांना लिहिलेले प्रसिद्ध “\emph{je le vois, mais je ne le crois
pas}” हे वाक्य त्याच्या संदर्भातून बरेचसे वेगळे केले जाते. तो संदर्भ पुढे दिला आहे (\citeauthor{Gouvea2011} यांच्याकडून उद्धृत):
\begin{quote}
या विषयाबद्दलच्या माझ्या उत्साहापोटी मी तुमच्या दयाळूपणाची आणि संयमाची इतकी मागणी करत असेन, तर मला क्षमा करा. मी तुम्हाला अलीकडे पाठवलेले विचार मलाही इतके अनपेक्षित, इतके नवे वाटतात की, आदरणीय मित्रा, त्यांची अचूकता तुम्ही निश्चित करेपर्यंत मला मनःशांती मिळू शकत नाही. तुम्ही माझ्याशी सहमत होत नाही तोपर्यंत मी एवढेच म्हणू शकतो:
\emph{je le vois, mais je ne le crois pas.}
\end{quote}
आपला निष्कर्ष “इतका अनपेक्षित, इतका नवा” आहे हे कँटर यांना माहीत होते. पण तो निष्कर्ष त्यांना खरोखर \emph{अविश्वसनीय} वाटत होता का, याबद्दल शंका आहे. \citeauthor{Gouvea2011} यांनी दाखवल्याप्रमाणे, त्यांनी डेडेकिंड यांना केवळ स्वतः दिलेल्या सिद्धतेची तपासणी करायला सांगितले होते.

भूमितीय अंतर्ज्ञानाच्या प्रश्नावर पाहिले तर, पियानो यांनी अवकाश भरणारा वक्र प्रकाशित करताना कोणतीही चित्रे दिली नव्हती. पण हिल्बर्ट यांनी स्वतःचा वक्र प्रकाशित केला तेव्हा आपला हेतू त्यांनी स्पष्ट केला: वाचकांनी “पुढील भूमितीय अंतर्ज्ञानाचा आधार घेतल्यास” पियानो यांचा निष्कर्ष त्यांना नीट समजावा, म्हणून त्यांनी \olref[his][set][pathology]{sec} मध्ये दिलेल्यांसारखी \emph{चित्रे} त्यात जोडली.

अधिक सर्वसाधारणपणे सांगायचे तर, प्रकाशित सिद्धतांमध्ये चित्रे वापरण्याची पद्धत काहीशी कमी झाली असली तरी गणितज्ञ सिद्धता करताना \emph{वारंवार} चित्रे वापरतात, ही गोष्ट नाकारता येत नाही. (ढोबळपणे सांगायचे तर, भूमितीय अंतर्ज्ञानावर कधी विसंबता येते हे चांगल्या गणितज्ञांना माहीत असते.)

थोडक्यात, अतिरंजित प्रचारावर विश्वास ठेवू नका; निदान तो विनातपासणी स्वीकारू नका. याबद्दल अधिक माहितीसाठी \citet{Giaquinto2007} वाचता येईल.

\end{document}
